% Reviewed translation: PDF pages 325--330(part) / printed pages 311--316(part).
\subsection{非奇异解分支的不可微逼近}
\label{subsec:navier-nondifferentiable-branches}

到目前为止, 我们一直假设逼近映射 $F_h$ 保留映射 $F$ 的光滑性,
因为逼近只作用于线性算子 $T$, 而不作用于 $G$. 但是有时必须用一个不再可微的映射
$G_h$ 逼近 $G$. 例如, 对 Navier--Stokes 方程的对流项采用迎风离散时就是如此.
下面的理论分析这种情形. 它是 Subsections~\ref{subsec:navier-nonsingular-branch-approximation}
和 \ref{subsec:navier-branch-nonlinear-application} 中理论的一个简单变体.
关于涵盖本节全部内容的更一般方法, 读者可参见 Rappaz~\cite{Rappaz1984}.

\MathBookSourcePage{326}
这里的情形比 Subsection~\ref{subsec:navier-nonsingular-branch-approximation}
稍为一般. 引入 $X$ 的闭子空间 $X_h$ 和 $\mathscr X$ 的闭子空间
$\mathscr X_h$, 并假设映射 $F_h$ 定义在 $\Lambda\times X_h$ 上、连续且取值于
$\mathscr X_h$. 对给定的 $\widetilde u_h\in X_h$ 和给定的
$\lambda\in\Lambda$, 以如下假设代替可微性:
\begin{equation}\label{eq:navier-branch-generalized-derivative-isomorphism}
\left\{
\begin{aligned}
&\text{there exists }\nabla_uF_h(\lambda,\widetilde u_h)
\in\mathcal L(X_h;\mathscr X_h)\\
&\text{which is an isomorphism from }X_h\text{ onto }\mathscr X_h.
\end{aligned}
\right.
\end{equation}
并且存在连续、单调递增的函数
$\mu\mapsto L_h(\lambda;\mu):\mathbb R_+\to\mathbb R_+$, 使
\begin{equation}\label{eq:navier-branch-generalized-lipschitz}
\begin{aligned}
&\|F_h(\lambda,v)-F_h(\lambda,w)
-\nabla_uF_h(\lambda,\widetilde u_h)\cdot(v-w)\|_{\mathscr X}\\
&\qquad\leq L_h(\lambda;\alpha)\|v-w\|_X
\quad\forall v,w\in S(\widetilde u_h;\alpha)\cap X_h.
\end{aligned}
\end{equation}
保留下列记号:
\begin{equation}\label{eq:navier-branch-generalized-local-quantities}
\left\{
\begin{aligned}
\varepsilon_h(\lambda)&=\|F_h(\lambda,\widetilde u_h)\|_{\mathscr X},\\
\gamma_h(\lambda)&=
\|[\nabla_uF_h(\lambda,\widetilde u_h)]^{-1}\|_{\mathcal L(\mathscr X_h;X_h)},
\end{aligned}
\right.
\end{equation}
其中
\[
\|B\|_{\mathcal L(\mathscr X_h;X_h)}
=\sup_{v\in\mathscr X_h}\frac{\|Bv\|_X}{\|v\|_{\mathscr X}}.
\]
于是有 Theorem~\ref{thm:navier-branch-local-existence-error} 的如下对应结果.

\setcounter{statement}{5}
\begin{Theorem}\label{thm:navier-branch-generalized-local-existence}
若 \eqref{eq:navier-branch-generalized-derivative-isomorphism} 和
\eqref{eq:navier-branch-generalized-lipschitz} 成立, 且另外有
\begin{equation}\tag{3.19$'$}\label{eq:navier-branch-generalized-contraction}
2\gamma_h(\lambda)L_h
(\lambda;2\gamma_h(\lambda)\varepsilon_h(\lambda))<1,
\end{equation}
则 Problem~\eqref{eq:navier-branch-discrete-problem} 有唯一解
$(\lambda,u_h(\lambda))$, 满足
\begin{equation}\tag{3.20$'$}\label{eq:navier-branch-generalized-local-ball}
u_h(\lambda)\in
S(\widetilde u_h;2\gamma_h(\lambda)\varepsilon_h(\lambda))\cap X_h.
\end{equation}
此外, 对每个满足 $\alpha\geq2\gamma_h(\lambda)\varepsilon_h(\lambda)$ 且
\begin{equation}\tag{3.22$'$}\label{eq:navier-branch-generalized-uniqueness-radius}
\gamma_h(\lambda)L_h(\lambda;\alpha)<1
\end{equation}
的 $\alpha$, $u_h(\lambda)$ 是较大球
$S(\widetilde u_h;\alpha)\cap X_h$ 中
Problem~\eqref{eq:navier-branch-discrete-problem} 的唯一解, 并有估计
\begin{equation}\tag{3.23$'$}\label{eq:navier-branch-generalized-error}
\|u_h(\lambda)-v_h\|_X
\leq\frac{\gamma_h(\lambda)}
{1-\gamma_h(\lambda)L_h(\lambda;\alpha)}
\|F_h(\lambda,v_h)\|_{\mathscr X}
\quad\forall v_h\in S(\widetilde u_h;\alpha)\cap X_h.
\end{equation}
\end{Theorem}

省略证明, 因为它与 Theorem~\ref{thm:navier-branch-local-existence-error}
的证明非常相似. 同样, 容易证明 Theorem~\ref{thm:navier-branch-uniform-branch}
的如下对应结果.

\setcounter{statement}{6}
\begin{Theorem}\label{thm:navier-branch-generalized-uniform-branch}
设 $\widetilde u_h$ 是给定的 $C^0(\Lambda;X_h)$ 函数, 满足: 对某个 $h_0>0$,
\begin{align}
\sup_{h\leq h_0}\left(\sup_{\lambda\in\Lambda}\gamma_h(\lambda)\right)
&=\gamma,
\tag{3.26$'$}\label{eq:navier-branch-generalized-uniform-gamma}\\
\lim_{h\to0}\left(\sup_{\lambda\in\Lambda}\varepsilon_h(\lambda)\right)&=0.
\tag{3.27$'$}\label{eq:navier-branch-generalized-uniform-residual}
\end{align}
若另外有
\setcounter{equation}{55}
\begin{equation}\label{eq:navier-branch-generalized-uniform-lipschitz}
\sup_{\lambda\in\Lambda}L_h(\lambda;\mu)=L_h(\mu)
\qquad\forall h\leq h_0,
\end{equation}
其中函数 $L_h(\mu)$ 关于 $\mu$ 和 $h$ 单调递增, 在 $\mu=0$ 处连续, 且满足
\begin{equation}\label{eq:navier-branch-generalized-zero-limit}
\lim_{h\to0}L_h(0)=0,
\end{equation}
则存在两个实常数 $\alpha>0$, $h_1>0$ 以及函数
$u_h\in C^0(\Lambda;X_h)$, 使对所有 $h\leq h_1$,
\[
\{(\lambda,u_h(\lambda));\lambda\in\Lambda\}
\text{ is a branch of solutions of }
\eqref{eq:navier-branch-discrete-problem}.
\]
对每个 $\lambda\in\Lambda$, $u_h(\lambda)$ 是
$S(\widetilde u_h(\lambda);\alpha)\cap X_h$ 中
Problem~\eqref{eq:navier-branch-discrete-problem} 的唯一解, 并有估计
\begin{equation}\tag{3.30$'$}\label{eq:navier-branch-generalized-uniform-error}
\|u_h(\lambda)-v_h\|_X
\leq2\gamma\|F_h(\lambda,v_h)\|_{\mathscr X}
\quad\forall v_h\in S(\widetilde u_h(\lambda);\alpha)\cap X_h.
\end{equation}
\end{Theorem}

\MathBookSourcePage{327}
\setcounter{statement}{7}
\begin{Remark}\label{rem:navier-branch-generalized-uniformity-in-h}
$L_h(\mu)$ 关于 $h$ 单调递增这一假设并非必要. 它可以替换为
$L_h(\mu)$ 在 $\mu=0$ 处关于 $h$ 一致连续这一条件.
\end{Remark}

为简单起见, 我们将用 Theorem~\ref{thm:navier-branch-generalized-uniform-branch}
求解比 \eqref{eq:navier-branch-operator-problem} 更窄的一类问题, 但读者容易把
下面的分析推广到一般情形. 更确切地说, 沿用
Subsection~\ref{subsec:navier-branch-nonlinear-application} 的记号, 我们要解
\begin{equation}\tag{3.33$'$}\label{eq:navier-branch-nondifferentiable-model}
F(\lambda,u)\equiv u+\lambda TG(u)=0,
\end{equation}
其中 $G$ 是从 $X$ 到 $Y$ 的 $C^1$ 映射,
$\lambda\in\Lambda$, $\Lambda$ 是 $\mathbb R$ 的紧区间. 算子 $T$ 保持不变,
但假设 $T\in\mathcal L(Y;V)$, 其中 $V$ 是 $X$ 的闭子空间.
这可以理解为对 $T$ 的一个正则性假设. 再假设问题
\[
F(\lambda,u(\lambda))=0
\]
有从 $\Lambda$ 到 $X$ 的非奇异解分支 $\lambda\mapsto u(\lambda)$, 且
$u\in C^0(\Lambda;V)$.

为进行逼近, 引入 $X$ 的闭子空间 $V_h$, 它赋予 $X$ 的范数;
并引入连续嵌入地包含 $Y$ 的空间 $Y_h$. 为避免混淆, 用
$\|\cdot\|_h$ 表示 $Y_h$ 的范数. 以从 $V_h$ 到 $Y_h$ 的 $C^0$ 映射
$G_h$ 逼近 $G$, 以算子 $T_h\in\mathcal L(Y_h;V_h)$ 逼近 $T$. 置
\setcounter{equation}{57}
\begin{equation}\label{eq:navier-branch-nondifferentiable-discrete-map}
F_h(\lambda,u_h)=u_h+\lambda T_hG_h(u_h),
\end{equation}
它是从 $\mathbb R\times V_h$ 到 $V_h$ 的 $C^0$ 映射.
下面的 Theorem 在不涉及 $G_h$ 可微性的情况下求解
Problem~\eqref{eq:navier-branch-nondifferentiable-model}.

\MathBookSourcePage{328}
\setcounter{statement}{7}
\begin{Theorem}\label{thm:navier-branch-nondifferentiable-approximation}
在下列假设下:

\textup{(i)}
\begin{align}
\|T_h\|_{\mathcal L(Y_h;V_h)}&\leq C,
\label{eq:navier-branch-th-bounded}\\
\lim_{h\to0}\|T-T_h\|_{\mathcal L(Y;X)}&=0.
\label{eq:navier-branch-th-converges}
\end{align}

\textup{(ii)} 存在算子 $\pi_h\in\mathcal L(V;V_h)$, 使
\begin{align}
\lim_{h\to0}\|v-\pi_hv\|_X&=0
&&\forall v\in V,
\label{eq:navier-branch-pi-converges}\\
\lim_{h\to0}\sup_{\lambda\in\Lambda}
\|G_h(\pi_hu(\lambda))-G(u(\lambda))\|_h&=0.
\label{eq:navier-branch-gh-consistency}
\end{align}

\textup{(iii)} 对所有 $u_h\in V_h$, 存在算子
$\nabla G_h(u_h)\in\mathcal L(V_h;Y_h)$, 使
\begin{align}
\lim_{h\to0}\sup_{\lambda\in\Lambda}
\|\nabla G_h(\pi_hu(\lambda))-DG(u(\lambda))\|_{\mathcal L(V_h;Y_h)}&=0,
\label{eq:navier-branch-generalized-gradient-consistency}\\
\|G_h(u_h)-G_h(u_h^*)-\nabla G_h(u_h^0)\cdot(u_h-u_h^*)\|_h
&\leq L_h(\mu;\|u_h^0\|_X)\|u_h-u_h^*\|_X
\label{eq:navier-branch-gh-generalized-lipschitz}
\end{align}
对所有 $u_h^0\in V_h$ 以及
$u_h,u_h^*\in S(u_h^0;\mu)\cap V_h$ 成立, 其中
$L_h:\mathbb R_+\times\mathbb R_+\to\mathbb R_+$ 关于每个变量以及关于 $h$
均为连续单调递增函数, 并满足
\begin{equation}\label{eq:navier-branch-lh-zero-limit}
\lim_{h\to0}L_h(0;\mu)=0
\qquad\forall\mu\in\mathbb R_+.
\end{equation}
则存在 $X$ 中原点的一个邻域 $\mathcal O$, 且对充分小的 $h\leq h_0$,
存在唯一的 $C^0$ 函数 $u_h:\Lambda\to V_h$, 使
\begin{equation}\label{eq:navier-branch-nondifferentiable-solution}
F_h(\lambda,u_h(\lambda))=0,
\qquad u_h(\lambda)\in u(\lambda)+\mathcal O
\quad\forall\lambda\in\Lambda.
\end{equation}
此外有误差估计
\begin{equation}\label{eq:navier-branch-nondifferentiable-error}
\begin{aligned}
\|u_h(\lambda)-u(\lambda)\|_X\leq C\big\{
&\|(I-\pi_h)u(\lambda)\|_X
+\|\lambda(T-T_h)G(u(\lambda))\|_X\\
&+\|G_h(\pi_hu(\lambda))-G(u(\lambda))\|_h\big\}.
\end{aligned}
\end{equation}
\end{Theorem}
\begin{Proof}
以 $\widetilde u_h=\widetilde u_h(\lambda)=\pi_hu(\lambda)$ 应用
Theorem~\ref{thm:navier-branch-generalized-uniform-branch}, 其中 $V_h$ 扮演 $X_h$
的角色. 因为假设 $u\in C^0(\Lambda;V)$, 所以
$\widetilde u_h\in C^0(\Lambda;V_h)$. 下面证明
\eqref{eq:navier-branch-generalized-uniform-residual}. 有
\[
\varepsilon_h(\lambda)=
\|F_h(\lambda,\widetilde u_h(\lambda))-F(\lambda,u(\lambda))\|_X,
\]
即
\begin{equation}\label{eq:navier-branch-nondifferentiable-residual-bound}
\begin{aligned}
\varepsilon_h(\lambda)leq
&\|\widetilde u_h(\lambda)-u(\lambda)\|_X
+|\lambda|\|T_h\{G_h(\widetilde u_h(\lambda))-G(u(\lambda))\}\|_X\\
&+|\lambda|\|(T_h-T)G(u(\lambda))\|_X.
\end{aligned}
\end{equation}

\MathBookSourcePage{329}
一方面, \eqref{eq:navier-branch-pi-converges} 和 Ascoli Lemma 蕴含
\[
\lim_{h\to0}\sup_{\lambda\in\Lambda}
\|\widetilde u_h(\lambda)-u(\lambda)\|_X=0.
\]
另一方面, \eqref{eq:navier-branch-th-bounded} 和
\eqref{eq:navier-branch-gh-consistency} 给出
\[
\lim_{h\to0}\sup_{\lambda\in\Lambda}
\|T_h\{G_h(\widetilde u_h(\lambda))-G(u(\lambda))\}\|_X=0.
\]
此外, 由映射 $G$ 的连续性和 \eqref{eq:navier-branch-th-converges} 可得
\[
\lim_{h\to0}\sup_{\lambda\in\Lambda}
\|(T_h-T)G(u(\lambda))\|_X=0.
\]
这三个极限立即给出 \eqref{eq:navier-branch-generalized-uniform-residual}.

现在验证 \eqref{eq:navier-branch-generalized-uniform-gamma}.
对每个 $u_h\in V_h$, 自然地定义
$\nabla_uF_h(\lambda,u_h)\in\mathcal L(V_h;V_h)$ 为
\begin{equation}\label{eq:navier-branch-generalized-f-gradient}
\nabla_uF_h(\lambda,u_h)=I+\lambda T_h\nabla G_h(u_h).
\end{equation}
于是可写成
\begin{equation}\label{eq:navier-branch-generalized-f-gradient-split}
\nabla_uF_h(\lambda,\widetilde u_h(\lambda))
=\lambda T_h\{\nabla G_h(\widetilde u_h(\lambda))-DG(u(\lambda))\}
+I+\lambda T_hDG(u(\lambda)).
\end{equation}
注意算子 $I+\lambda T_hDG(u(\lambda))$ 同时属于
$\mathcal L(X;X)$ 和 $\mathcal L(V_h;V_h)$. 此外,
\[
I+\lambda T_hDG(u(\lambda))-D_uF(\lambda,u(\lambda))
=\lambda(T_h-T)DG(u(\lambda)),
\]
由 \eqref{eq:navier-branch-th-converges} 和 $G$ 的可微性,
\[
\lim_{h\to0}\sup_{\lambda\in\Lambda}
\|(T_h-T)DG(u(\lambda))\|_{\mathcal L(X;X)}=0.
\]
所以应用 Lemma~\ref{lem:navier-branch-derivative-perturbation} 可知,
存在实数 $h_0>0$, 使对所有 $h\leq h_0$ 及所有 $\lambda\in\Lambda$,
$I+\lambda T_hDG(u(\lambda))$ 既是 $X$ 上的同构, 也是 $V_h$ 上的同构. 并且
\begin{equation}\label{eq:navier-branch-restricted-inverse-bound}
\begin{aligned}
\|(I+\lambda T_hDG(u(\lambda)))^{-1}\|_{\mathcal L(V_h;V_h)}
&\leq
\|(I+\lambda T_hDG(u(\lambda)))^{-1}\|_{\mathcal L(X;X)}\\
&\leq2\gamma(\lambda)
\qquad\forall h\leq h_0,\quad\forall\lambda\in\Lambda,
\end{aligned}
\end{equation}
其中 $\gamma(\lambda)$ 由 \eqref{eq:navier-branch-gamma} 定义.
同样, 由 \eqref{eq:navier-branch-th-bounded} 和
\eqref{eq:navier-branch-generalized-gradient-consistency} 可得
\[
\lim_{h\to0}\sup_{\lambda\in\Lambda}
\|T_h\{\nabla G_h(\widetilde u_h(\lambda))-DG(u(\lambda))\}
\|_{\mathcal L(V_h;V_h)}=0.
\]
再次应用 Lemma~\ref{lem:navier-branch-derivative-perturbation}, 并使用
\eqref{eq:navier-branch-generalized-f-gradient-split} 和
\eqref{eq:navier-branch-restricted-inverse-bound}, 可知存在另一个满足
$0<h_1\leq h_0$ 的 $h_1$, 使对所有 $h\leq h_1$ 和
$\lambda\in\Lambda$, $\nabla_uF_h(\lambda,\widetilde u_h(\lambda))$
是 $V_h$ 上的同构, 且
\[
\|[\nabla_uF_h(\lambda,\widetilde u_h(\lambda))]^{-1}\|_{\mathcal L(V_h;V_h)}
\leq4\gamma(\lambda).
\]
因 $\gamma(\lambda)$ 有界, 这就证明了
\eqref{eq:navier-branch-generalized-uniform-gamma}.

\MathBookSourcePage{330}
还需验证由 \eqref{eq:navier-branch-generalized-lipschitz} 定义的函数 $L_h$
所满足的条件. 由 \eqref{eq:navier-branch-th-bounded} 和
\eqref{eq:navier-branch-gh-generalized-lipschitz} 有上界
\begin{align*}
&\|\nabla_uF_h(\lambda,\widetilde u_h(\lambda))\cdot(v_h-w_h)
-F_h(\lambda,v_h)+F_h(\lambda,w_h)\|_X\\
&\qquad\leq C|\lambda|L_h(\alpha;\|\widetilde u_h(\lambda)\|_X)
\|v_h-w_h\|_X,
\end{align*}
对所有 $\lambda\in\Lambda$ 以及所有
$v_h,w_h\in S(\widetilde u_h(\lambda);\alpha)\cap V_h$ 成立. 令
\[
K=\sup_{\substack{h\leq h_1\\\lambda\in\Lambda}}
\|\widetilde u_h(\lambda)\|_X.
\]
$L_h$ 的单调性蕴含
\[
L_h(\alpha;\|\widetilde u_h(\lambda)\|_X)\leq L_h(\alpha;K).
\]
由假设, 映射 $\mu\mapsto L_h(\mu;K)$ 关于 $h$ 和 $\mu$ 单调递增,
并在 $\mu=0$ 处连续. 此外, 由 \eqref{eq:navier-branch-lh-zero-limit},
它满足 \eqref{eq:navier-branch-generalized-zero-limit}.

因此可应用 Theorem~\ref{thm:navier-branch-generalized-uniform-branch} 的结论.
它给出 \eqref{eq:navier-branch-nondifferentiable-solution}; 随后
\eqref{eq:navier-branch-nondifferentiable-error} 由
\eqref{eq:navier-branch-generalized-uniform-error} 取
$v_h=\widetilde u_h(\lambda)$, 再结合
\eqref{eq:navier-branch-nondifferentiable-residual-bound} 和
\eqref{eq:navier-branch-th-bounded} 得到.
\end{Proof}

\setcounter{statement}{8}
\begin{Remark}\label{rem:navier-branch-regularity-space-v}
空间 $V$ 只出现在假设 \eqref{eq:navier-branch-pi-converges} 中,
因而在这个证明中仅起次要作用. 当限制算子 $\pi_h$ 要求的正则性高于空间 $X$
所能提供的正则性时, 它是有用的.
\end{Remark}
